\chapter{مرور پیشینه پژوهش}
\label{ch:literature}

\section{مقدمه}

در این فصل، پژوهش‌های مرتبط با حل مسائل ترکیبیاتی با استفاده از یادگیری تقویتی، شبکه‌های عصبی گرافی و همچنین روش‌های موجود برای مسئله مجموعه احاطه‌گر و مجموعه احاطه‌گر کامل مرور می‌شوند.

\section{حل مسائل ترکیبیاتی با یادگیری تقویتی و شبکه‌های عصبی گرافی}

استفاده از یادگیری عمیق برای حل مسائل بهینه‌سازی ترکیبیاتی در سال‌های اخیر مورد توجه زیادی قرار گرفته است. خلیل و همکاران \cite{khalil2017learning} رویکردی مبتنی بر $DQN$ و $GNN$ برای حل مسائل NP-سخت مانند مجموعه مستقل بیشینه، حداقل پوشش رأسی و فروشنده دوره‌گرد ارائه دادند.

بللو و همکاران \cite{bello2016neural} با استفاده از یادگیری تقویتی و رمزگذار-رمزگشا\footnote{Encoder-Decoder}، روشی برای بهبود راه‌حل‌های مسئله فروشنده دوره‌گرد معرفی کردند. وینیالز و همکاران \cite{vinyals2015pointer} نیز مدل $Pointer\ Network$ را برای تولید دنباله‌های مرتب‌شده در مسائل ترکیبیاتی پیشنهاد دادند.

همچنین بنگیو و همکاران \cite{bengio2021} مروری جامع بر کاربرد یادگیری ماشین در بهینه‌سازی ترکیبیاتی ارائه نمودند و نشان دادند که روش‌های یادگیری‌محور می‌توانند در مواردی از روش‌های سنتی پیشی بگیرند یا مکمل آن‌ها باشند.

\section{پژوهش‌های مرتبط با مسئله مجموعه احاطه‌گر}

مسئله مجموعه احاطه‌گر از دیرباز مورد توجه پژوهشگران بوده است. روش‌های حریصانه، الگوریتم‌های تقریبی و فراابتکاری متعددی برای این مسئله ارائه شده‌اند \cite{haynes1998fundamentals}.

در سال‌های اخیر، رویکردهای یادگیری‌محور نیز ظهور کرده‌اند. چن و همکاران \cite{chen2024dsp} اولین بار چارچوبی یکپارچه مبتنی بر شبکه عصبی گرافی و $DDQN$ را برای حل مسئله مجموعه احاطه‌گر ارائه دادند. در این روش، مسئله به‌صورت فرآیند تصمیم‌گیری ترتیبی مدلسازی شده و عامل یاد می‌گیرد با انتخاب رأس‌های مناسب، مجموعه احاطه‌گر با اندازه کوچک تولید کند.

این روش در مقایسه با الگوریتم‌های حریصانه و برخی روش‌های فراابتکاری، نتایج قابل‌قبولی روی خانواده‌های مختلف گراف از جمله گراف‌های تصادفی $ER$، درخت‌ها و شبکه‌های واقعی نشان داده است.

\section{پژوهش‌های مرتبط با مسئله مجموعه احاطه‌گر کامل}

مسئله مجموعه احاطه‌گر کامل نیز موضوع پژوهش‌های فراوانی بوده است. تحقیقات زیادی بر روی کران‌های بالا و پایین، پیچیدگی محاسباتی و الگوریتم‌های تقریبی این مسئله متمرکز بوده‌اند \cite{haynes1998fundamentals, cockayne2004total}. با این حال، تاکنون تلاش محدودی برای حل این مسئله با استفاده از یادگیری تقویتی و شبکه‌های عصبی گرافی انجام شده است.

با توجه به شباهت ساختاری بین مسئله مجموعه احاطه‌گر و مجموعه احاطه‌گر کامل، طبیعی است که بررسی شود آیا چارچوب پیشنهادی در \cite{chen2024dsp} قابل گسترش به نسخه کامل این مسئله باشد یا خیر. این پرسش، انگیزه اصلی پژوهش حاضر را تشکیل می‌دهد.

\section{خلاصه}

در این فصل، پیشینه پژوهش در سه حوزه مرور شد: حل مسائل ترکیبیاتی با RL و GNN، روش‌های موجود برای مسئله مجموعه احاطه‌گر و روش‌های مرتبط با مجموعه احاطه‌گر کامل. بررسی‌ها نشان می‌دهد که اگرچه روش‌های یادگیری‌محور برای DSP موفق بوده‌اند، گسترش مستقیم آن‌ها به TDSP نیازمند بازنگری در فضای عمل، حالت و تابع پاداش است که در فصل \ref{ch:tdsp-extension} به آن پرداخته می‌شود.
